\uuid{5f8n}
\titre{Équation droite/plan}
\auteur{bodin ; hanani ; mzari}
\organisation{exo7}
\licence{CC BY-NC-SA 4.0}
\datecreate{2020-01-20}
\niveau{L1}
\module{Géométrie}
\chapitre{Géométrie de l'espace}
\sousChapitre{Plans}
\difficulte{3}
\theme{Plans -- Droites, équation droite, plan}
\exo7qcmid{lille-1.geometrie-espace.019}

\contenu{
\texte{
Pour ces questions, l'espace est muni d'un repère orthonormé direct  $(O, \vec{i}, \vec{j}, \vec{k})$.
}

\begin{question}
\qclassification{141.03,141.04}
\qtags{motcle=équation droite/plan}
On considère les droites $D_1 : \left\{\begin{array}{ccl}x&=&z-1\\y&=&2z+1 \end{array}\right.$ et
$D_2 : \left\{\begin{array}{ccl}y&=&3x\\z&=&1 \end{array}\right.$. Soit $P_1$ et $P_2$ des plans parallèles contenant $D_1$ et $D_2$ respectivement.    Quelles sont les assertions vraies ?
\begin{answers}

       \good{Une équation cartésienne de $P_1$ est $ 3x-y-z+4=0$.}

      \bad{Une équation cartésienne de $P_1$ est $ 4x-y-z+5=0$.}

     \bad{Une équation cartésienne de $P_2$ est $ 4x-y-z+1=0$.}

     \good{Une équation cartésienne de $P_2$ est $ 3x-y-z+1=0$.}

\end{answers}
\begin{explanations}
$D_1$ passe par le point $A_1(-1,1,0)$ et est dirigée par le vecteur $\vec{u_1}(1,2,1)$.
$D_2$ passe par le point $A_2(0,0,1)$ et est dirigée par le vecteur $\vec{u_2}(1,3,0)$. $P_1$ passe donc par $A_1$ est de vecteur normal  $\vec{n}=\vec{u_1} \wedge \vec{u_2}$ et $P_2$ passe donc par $A_2$ est de vecteur normal  $\vec{n}$.
\end{explanations}
\end{question}
}
